\documentclass[twocolumn]{aastex631}
\begin{document}
\title{The reionization ionizing-photon-budget shortfall for star-forming galaxies is not robust to systematics at z\~{}5 (lowxi systematic corner)}
\author{NebulaMind Lab (autonomous pipeline)}
\affiliation{NebulaMind Lab, \url{https://lab.nebulamind.net}}
\begin{abstract}
Reionization ionizing-photon-budget reconciliation at z\~{}5: star-forming galaxies require f\_esc=0.039 (+0.039/-0.020) to close the budget (Madau-Dickinson SFRD, log xi\_ion=25.5+/-0.15, clumping C=2-5, JWST-SFRD tail), versus indirect-proxy-inferred f\_esc=0.062 (+0.108/-0.039) from LzLCS O32/beta calibrations. Median delta(required-inferred)=-0.021 dex-frac (16-84\%: -0.128 to +0.029); 35\% of the systematic MC shows a shortfall. Conclusion: the budget CLOSES within the systematic, and the sign holds under both O32 and beta calibrations. The result is bounded by the xi\_ion x clumping x proxy-calibration systematic, not by statistics. Generated autonomously from public data (jwst) via the NebulaMind Lab runner: a bounded, reproducible, descriptive study.
\end{abstract}
\section{Introduction}
Recent studies have highlighted potential discrepancies in the reionization photon budget, suggesting that star-forming galaxies may not produce enough ionizing photons to account for the observed reionization [Muñoz2024]. This has led to a renewed interest in understanding the role of these galaxies in the reionization process. Previous work has focused on calibrating excursion set reionization models to conserve ionizing photons [Park2022] and assessing the galaxy ionizing photon budget at high redshifts [Duncan2015]. However, uncertainties remain regarding the escape fraction of ionizing photons from these galaxies.
\section{Data and method}
To address this issue, we adopt a literature-anchored budget calculation approach, relying on published values for key parameters. The cosmic star formation rate density (SFRD) is taken from the analytic fitting function provided by Madau \& Dickinson (2014). We use published calibrations for the ionizing photon production efficiency (xi\_ion) and the escape fraction (f\_esc) proxy, based on O32/beta ratios [LzLCS: Chisholm+22, Flury+22; Simmonds+24]. This method allows us to reconcile the reionization photon budget without relying on new observational data.
\section{Result}
Our analysis reveals that star-forming galaxies at z\~{}5 require an escape fraction of f\_esc=0.039 (+0.039/-0.020) to close the ionizing-photon-budget, assuming a Madau-Dickinson SFRD and log xi\_ion=25.5±0.15, with clumping factors C between 2-5. This value is compared to indirect-proxy-inferred f\_esc=0.062 (+0.108/-0.039) derived from LzLCS O32/beta calibrations. The median difference between the required and inferred escape fractions is -0.021 dex-frac, with a range of -0.128 to +0.029 across 16-84\% of systematic Monte Carlo realizations. Notably, 35\% of these simulations show a shortfall in the photon budget.
\begin{figure}\centering\includegraphics[width=\columnwidth]{result.png}\caption{The reionization ionizing-photon-budget shortfall for star-forming galaxies is not robust to systematics at z\~{}5 (lowxi systematic corner)}\end{figure}
\section{Caveats}
It is essential to acknowledge the limitations of our approach. The reliance on published calibrations and assumptions about xi\_ion and clumping factors introduces uncertainties that are not fully quantifiable. Additionally, the use of a single selection criterion for star-forming galaxies may introduce biases, as it does not account for variations in galaxy properties or environments. Furthermore, the lack of direct observational data to validate our results means that they remain dependent on the accuracy of previous studies. These caveats highlight the need for further research and more robust measurements to confirm our findings.
\section*{References}
\footnotesize
[Muoz2024] Muñoz et al. (2024). Reionization after JWST: a photon budget crisis?. bibcode 2024MNRAS.535L..37M \par [Park2022] Park et al. (2022). Calibrating excursion set reionization models to approximately conserve ionizing photons. bibcode 2022MNRAS.517..192P \par [Duncan2015] Duncan et al. (2015). Powering reionization: assessing the galaxy ionizing photon budget at z < 10. bibcode 2015MNRAS.451.2030D \par [Davies2021] Davies et al. (2021). The Predicament of Absorption-dominated Reionization: Increased Demands on Ionizing Sources. bibcode 2021ApJ...918L..35D \par [Madau2017] Madau (2017). Cosmic Reionization after Planck and before JWST: An Analytic Approach. bibcode 2017ApJ...851...50M

\end{document}
